\subsubsection{Discrete Abel--Weil Explicit Formula and Polar Rigidity: Characterizing RH via Prime-End Generating Functions}\label{subsec:discrete-abel-weil-polar-rigidity}
This section presents a collection of results that can be verified line by line: starting from the Chebyshev function on the prime end, via ``logarithmic sampling + discrete Abel transform'' one constructs an analytic/meromorphic generating function on the unit disk, whose pole positions in the \(r\)-plane are in strict one-to-one correspondence with the zeros of \(\zeta\), thereby yielding a ``unit disk holomorphy \(\Leftrightarrow\) RH'' equivalence criterion and casting the ``scan readout'' as a Hardy boundary problem and discrete spectral inversion framework.

\paragraph{Chebyshev function and logarithmic sampling error sequence}
Denote the von Mangoldt function by \(\Lambda(n)\) and define
\[
\psi(x):=\sum_{n\le x}\Lambda(n).
\]
To avoid technical issues at jump discontinuities, take the standard half-weight version \(\psi_0(x)\) (averaging left and right limits at prime powers); its difference from \(\psi(x)\) occurs only on a discrete point set and does not affect the ``analytic domain/pole position'' conclusions below.
Fix an arbitrary base \(b>1\). Define the logarithmically sampled and \(1/2\)-normalized error sequence
\[
\boxed{\ E_b(n):=b^{-n/2}\bigl(\psi_0(b^n)-b^n\bigr),\qquad n=0,1,2,\dots\ }
\]
together with its Abel generating function
\[
\boxed{\ \mathcal A_b(r):=\sum_{n\ge 0}E_b(n)\,r^n.\ }
\]

\paragraph{Minimal-input form of the explicit formula}
This section uses only the following standard decomposition of the explicit formula (available in any analytic number theory textbook; only its structure is needed here, not the sharpest error term):
there exists a function \(G(x)\) composed of trivial zeros and a constant term (e.g., one may take \(G(x)=-\log(2\pi)-\tfrac12\log(1-x^{-2})\) in a suitably interpreted version) such that for all \(x>1\),
\begin{equation}\label{eq:psi0-explicit-schematic}
\boxed{\ \psi_0(x)=x-\sum_{\rho}\frac{x^\rho}{\rho}+G(x)\ },
\end{equation}
where \(\rho\) runs over all nontrivial zeros of \(\zeta(s)\), counted with multiplicity and paired with conjugates. Moreover, \(G(x)=O(1)\) (more precisely, \(G(x)=-\log(2\pi)+O(x^{-2})\)).

\begin{theorem}[Discrete Abel--Weil expansion: zeros \(\to\) \(r\)-plane poles]\label{thm:discrete-abel-weil-expansion}
In the region where \(|r|\) is sufficiently small for the exchange of summation to be justified, one has the decomposition
\begin{equation}\label{eq:Abel-Weil-decomposition}
\boxed{\ \mathcal A_b(r)
=-\sum_{\rho}\frac{1}{\rho}\cdot\frac{1}{1-r\,b^{\rho-\frac12}}
\;+\;H_b(r)\ },
\end{equation}
where the remainder
\[
H_b(r):=\sum_{n\ge 0}b^{-n/2}G(b^n)\,r^n
\]
is analytic at least in the disk \(|r|<b^{1/2}\) (in fact its power series convergence radius is \(\ge b^{1/2}\)).
Therefore each nontrivial zero \(\rho\) induces a pole orbit
\[
1-r\,b^{\rho-\frac12}=0
\quad\Longleftrightarrow\quad
\boxed{\ r_\rho=b^{-(\rho-\frac12)}\ },
\]
with modulus
\[
\boxed{\ |r_\rho|=b^{-(\Re\rho-\frac12)}\ }.
\]
\end{theorem}
\begin{proof}
From \eqref{eq:psi0-explicit-schematic} with \(x=b^n\), multiplied by \(b^{-n/2}\), one obtains
\[
E_b(n)
=-\sum_\rho \frac{b^{(\rho-\frac12)n}}{\rho}
\;+\;b^{-n/2}G(b^n).
\]
In the disk \(|r|\) sufficiently small, the summation over \(n\) and \(\rho\) may be exchanged via locally uniform convergence (first exchange for a fixed truncation of the zero set, then remove the truncation using zero-counting and summation by parts; only locally uniform convergence on \(|r|<R\) is needed). Thus
\[
\sum_{n\ge 0}E_b(n)r^n
=-\sum_\rho \frac{1}{\rho}\sum_{n\ge 0}\bigl(r\,b^{\rho-\frac12}\bigr)^n
\;+\;\sum_{n\ge 0}b^{-n/2}G(b^n)r^n.
\]
The inner sum in the first term is a geometric series \(\sum_{n\ge 0}(r\,b^{\rho-\frac12})^n=1/(1-r\,b^{\rho-\frac12})\), yielding \eqref{eq:Abel-Weil-decomposition}.
For the remainder, since \(G(b^n)=O(1)\), one has \(b^{-n/2}G(b^n)=O(b^{-n/2})\), so \(H_b(r)\) is analytic in \(|r|<b^{1/2}\).
\end{proof}

\paragraph{Polar rigidity and the geometric criterion for RH (unit disk holomorphy)}
Let
\[
\sigma_*:=\sup_{\rho}\Re\rho
\]
denote the supremum of the real parts of nontrivial zeros (RH is equivalent to \(\sigma_*=\tfrac12\)).

\begin{theorem}[Analytic radius captures the rightmost zero real part]\label{thm:radius-captures-sigma-star}
Let \(R_b\) denote the maximal radius of the ``pole-free disk'' centered at \(r=0\) for the power series expansion of \(\mathcal A_b(r)\). Then
\[
\boxed{\ R_b=b^{-(\sigma_*-\frac12)}\ }.
\]
More specifically:
\begin{itemize}
\item If \(|r|<R_b\), then no pole induced by nontrivial zeros lies in this disk, so \(\mathcal A_b(r)\) is holomorphic on \(|r|<R_b\);
\item For any \(\varepsilon>0\), in the larger disk \(|r|<R_b+\varepsilon\), one generally encounters poles induced by the ``rightmost zeros'' (at least in the meromorphic sense, holomorphy fails).
\end{itemize}
In particular,
\[
\boxed{\ \mathrm{RH}\ \Longleftrightarrow\ \mathcal A_b(r)\ \text{is holomorphic on}\ |r|<1.\ }
\]
The quantified version is: for any \(\delta\ge 0\),
\[
\boxed{\ \bigl(\forall\rho,\ \Re\rho\le\tfrac12+\delta\bigr)
\Longleftrightarrow
\mathcal A_b(r)\ \text{is holomorphic on}\ |r|<b^{-\delta}.\ }
\]
\end{theorem}
\begin{proof}
By Theorem \ref{thm:discrete-abel-weil-expansion}, each nontrivial zero \(\rho\) induces a pole \(r_\rho=b^{-(\rho-\frac12)}\), so the set of all pole-to-origin distances is \(\{|r_\rho|\}_\rho=\{b^{-(\Re\rho-\frac12)}\}_\rho\). The maximal ``pole-free disk'' radius for the holomorphic continuation from the center \(0\) equals the infimum of all pole moduli, i.e.,
\[
R_b=\inf_\rho |r_\rho|=\inf_\rho b^{-(\Re\rho-\frac12)}=b^{-(\sup_\rho \Re\rho-\frac12)}=b^{-(\sigma_*-\frac12)}.
\]
The quantified version follows by the same argument.
\end{proof}

\paragraph{Hardy boundary problem and the regularization framework for ``unitary scan''}
Under RH, the poles of \(\mathcal A_b\) lie on the boundary \(|r|=1\), making ``boundary integrability/finite energy'' require regularization. To this end, introduce an exponential damping parameter \(\delta>0\) that pushes the unit circle boundary singularities outside the disk.

\begin{definition}[Damped Abel generating function]\label{def:Abel-damped}
For any \(\delta\ge 0\), define
\[
\boxed{\ \mathcal A_{b,\delta}(r):=\mathcal A_b(b^{-\delta}r)=\sum_{n\ge 0}E_b(n)\,b^{-\delta n}\,r^n.\ }
\]
\end{definition}

\begin{proposition}[Pole positions after damping]\label{prop:damped-poles}
One has the decomposition
\[
\mathcal A_{b,\delta}(r)
=-\sum_{\rho}\frac{1}{\rho}\cdot\frac{1}{1-r\,b^{\rho-(\frac12+\delta)}}\;+\;H_{b,\delta}(r),
\]
where \(H_{b,\delta}(r)=H_b(b^{-\delta}r)\) is analytic at least on \(|r|<b^{1/2+\delta}\). The corresponding poles are
\[
\boxed{\ r_{\rho,\delta}=b^{\delta}r_\rho=b^{-(\rho-(\frac12+\delta))}\ },\qquad
|r_{\rho,\delta}|=b^{-(\Re\rho-(\frac12+\delta))}.
\]
Therefore
\[
\bigl(\forall\rho,\ \Re\rho\le\tfrac12+\delta\bigr)
\Longleftrightarrow
\mathcal A_{b,\delta}\ \text{is holomorphic on}\ |r|<1.
\]
\end{proposition}

\begin{theorem}[Hardy energy criterion \(\Leftrightarrow\) strict surplus beyond the rightmost zero]\label{thm:hardy-criterion-delta}
Fix \(\delta\ge 0\). The following two conditions are strictly equivalent:
\begin{enumerate}
\item There exists \(\varepsilon>0\) such that \(\forall\rho,\ \Re\rho\le \tfrac12+\delta-\varepsilon\);
\item \(\mathcal A_{b,\delta}\) is holomorphic on some disk \(|r|<1+\eta\) (\(\eta>0\)) (equivalently, it has no pole near the unit circle), so that \(\mathcal A_{b,\delta}\in H^2(\mathbb D)\) and its radial limit belongs to \(L^2\) on \(\TT\).
\end{enumerate}
In particular, RH can be equivalently restated as an ``arbitrarily small damping Hardy finite energy'' criterion:
\[
\boxed{\ \mathrm{RH}\ \Longleftrightarrow\ \forall\delta>0,\ \mathcal A_{b,\delta}\in H^2(\mathbb D).\ }
\]
\end{theorem}
\begin{proof}
\((1)\Rightarrow(2)\): By Proposition \ref{prop:damped-poles}, the pole moduli satisfy
\[
|r_{\rho,\delta}|=b^{-(\Re\rho-(\frac12+\delta))}\ge b^{\varepsilon}>1,
\]
so there exists \(\eta>0\) such that all poles lie outside \(|r|\ge 1+\eta\), and thus \(\mathcal A_{b,\delta}\) is holomorphic on \(|r|<1+\eta\). By Cauchy estimates for analytic functions on the larger disk, its Taylor coefficients decay geometrically, whence \(\sum_{n\ge 0}|E_b(n)b^{-\delta n}|^2<\infty\), equivalent to \(\mathcal A_{b,\delta}\in H^2(\mathbb D)\).

\((2)\Rightarrow(1)\): If some zero satisfies \(\Re\rho\ge \tfrac12+\delta\), then by Proposition \ref{prop:damped-poles}, \(\mathcal A_{b,\delta}\) has a pole within the closed disk \(|r|\le 1\) (inside the open disk when \(\Re\rho>\tfrac12+\delta\); on the boundary when equality holds). However, an \(H^2(\mathbb D)\) function must be holomorphic on the open disk, and its \(L^2\) boundary values exclude first-order polar singularities on the unit circle (otherwise \(\sup_{0<r<1}\int_0^{2\pi}|f(re^{i\theta})|^2\,\dd\theta=\infty\)). Therefore \(\mathcal A_{b,\delta}\in H^2\) implies no zero satisfies \(\Re\rho\ge \tfrac12+\delta\); by compactness of the surplus, some \(\varepsilon>0\) exists such that \(\Re\rho\le\tfrac12+\delta-\varepsilon\).

The final ``RH \(\Leftrightarrow\) \(\forall\delta>0\), \(H^2\) regularity'' statement follows from \((1)\Leftrightarrow(2)\) by taking \(\delta>0\) arbitrarily small and setting \(\varepsilon=\delta/2\).
\end{proof}

\paragraph{Boundary angular frequency and imaginary-part encoding of zeros (scan readout)}
Viewing \(\mathcal A_b(r)\) as a function first defined on \(|r|<1\), then taking the radial limit \(r=e^{i\theta}\) (understanding boundary singularities in the distributional sense): if RH holds and \(\rho=\tfrac12+i\gamma\), then
\[
b^{\rho-\frac12}=b^{i\gamma}=e^{i\gamma\log b},
\qquad
r_\rho=e^{-i\gamma\log b}\in\TT,
\]
and the pole angular position satisfies
\[
\boxed{\ \theta\equiv -\gamma\log b\pmod{2\pi}\ }.
\]
This rigorously encodes the ``zero imaginary part \(\gamma\)'' as an angular frequency on the unit circle boundary, thereby casting the ``scan readout'' as a Hardy-type boundary spectral problem.

\paragraph{Discrete spectral inversion: approximate recovery of zeros from prime block data}
From the coefficient representation of Theorem \ref{thm:discrete-abel-weil-expansion} one sees (under RH) that
\[
E_b(n)=-\sum_{\rho}\frac{b^{i\gamma n}}{\rho}+O(b^{-n/2}),
\]
i.e., \(E_b(n)\) is a sum of complex exponentials with a geometrically decaying tail. Therefore, to recover zeros with \(|\gamma|\le T\), it suffices to take the sampling length \(N\) at the Nyquist scale
\[
N\ \gtrsim\ \frac{T\log b}{2\pi}
\]
and use \(\{E_b(0),\dots,E_b(N-1)\}\) as input to standard frequency estimation tools (Prony, matrix pencil, ESPRIT, etc.) to approximately recover the frequency set \(\{\gamma\log b\}\), then divide by \(\log b\) to obtain estimates of \(\gamma\). The remainder \(O(b^{-n/2})\) provides a natural noise floor that decays geometrically with \(n\).

\endinput
